% !TeX root = ../main.tex

\xchapter{结论与展望}{Conclusion and Future Work}

\xsection{结论}{Conclusion}

随着深度学习模型规模与算力需求持续增长，多核神经网络芯片已成为大模型训练与推理的主流算力载体，但其在流片前的端到端性能评估始终是一道难题。硬件仿真平台（如 Synopsys ZeBu）精度高但成本昂贵、配置僵化，难以承担大规模设计空间探索。纯解析模型又无法刻画核间通信、片外存储与多核同步带来的系统级动态行为。因此，构建一套既能承接调度器输出、又能在多档精度间切换、并支持系统级瓶颈归因的周期级模拟器，成为本研究的出发点。

围绕这一目标，本文面向多核神经网络芯片设计并实现了一套以调度器中间表示驱动的模块化、可插拔的周期级模拟器。模拟器以编译/调度器输出的统一任务表示为唯一输入，避免了踪迹驱动方法在架构调整后反复重生成踪迹的代价。模拟器通过抽象接口将片上网络与片外存储解耦为简化与高保真两档后端，并以五态状态机驱动核内加载、计算和写回流水。另外，模拟器还以引用计数管理片上 SRAM 生命周期，以全局事件驱动与多速率推进处理跨时钟域同步，从而在多核并发与资源争用下保持因果一致与可复现性。在模拟器基础之上，本文进一步完成了从精度验证到真实负载瓶颈分解，再到多核递进式优化与三维设计空间探索的完整实验链路。

实验结果表明，本文所构建的模拟器在精度、可解释性与设计空间探索能力三方面均得到了系统性验证。在精度上，与 Synopsys ZeBu 硬件模拟平台在相同调度与配置下的总周期对照表明，本模拟器在 ResNet34、ResNet50 等真实负载下的相对误差控制在较小范围内。另外，简化后端与高保真后端形成了稳定的精度梯度，满足从快速校验到精细评估的不同需求。在瓶颈归因上，端到端阶段分解与各核热力图清晰刻画了 ResNet 类负载中计算与访存的相对占比，并揭示出 ResNet50 相对于 ResNet34 因层数加深而呈现出的访存压力上升趋势。在系统级优化与设计空间探索上，模拟器在四核测试芯片场景下通过片外 DRAM 带宽提升到 SRAM 端口加宽的递进式优化，将 ResNet50 总时钟周期从 238\,803 降至 62\,818，取得 3.80 倍加速，并清晰呈现出瓶颈从片外带宽逐级迁移至片上供数的过程。沿 MAC 规模、核数与批量大小三个维度的 12 组配置探索则进一步指出，在存储受限状态下提升算力对总时钟周期无明显改善。另外，核数扩展存在由负载并行度决定的上限，批量扩展效率仅约 24\%。这些结论不仅验证了模拟器在流片前架构定界中的实际价值，也为多核神经网络芯片在算力、互连与存储之间的权衡提供了可操作的设计依据。

\xsection{展望}{Future Work}

随着大模型对算力需求的不断增长，多核神经网络芯片的设计空间也变得越来越复杂。如何在流片前对这类芯片做出更可靠的评估和取舍，仍然是一个值得长期关注的问题。未来，相关研究有望在体系结构、编译调度和负载形态等方面进一步融合，让设计、评估和优化这几个环节的配合更加顺畅。本文的工作只是这个方向上的一次尝试，希望能为后续的研究和工程实践提供一些参考。
